\documentclass[11pt,a4paper]{article}

% JITE-oriented manuscript settings
\usepackage[a4paper,top=3.3cm,bottom=3.3cm,left=3.5cm,right=3.5cm]{geometry}
\usepackage[T1]{fontenc}
\usepackage[utf8]{inputenc}
\usepackage{lmodern}
\usepackage{setspace}
\singlespacing
\usepackage{amsmath,amssymb,amsthm}
\usepackage[backend=biber,style=authoryear,natbib=true,sorting=nyt,maxbibnames=99,giveninits=true,uniquename=init,doi=false,isbn=false,url=false,eprint=false]{biblatex}
\usepackage{hyperref}
\usepackage{enumitem}
\usepackage{microtype}
\usepackage{url}

\hypersetup{hidelinks}

\setlength{\parindent}{1.5em}
\setlength{\parskip}{0.35em}
\setlist[enumerate]{itemsep=0.2em,topsep=0.3em}

\newtheorem{proposition}{Proposition}
\newtheorem{remark}{Remark}

\title{External Agentic-AI Platforms, Retained Human Recovery Capability, and the Allocation of Authority}
\author{Ryota Matsuki\thanks{Independent Researcher, Matsuyama, Ehime, Japan. Corresponding author. E-mail: \href{mailto:ryota.matsuki@gmail.com}{ryota.matsuki@gmail.com}. ORCiD: \url{https://orcid.org/0009-0005-2329-531X}.}}
\date{Submission date: March 17, 2026}

\addbibresource{refs.bib}

\begin{document}

\maketitle

\begin{abstract}
This paper studies how firms allocate authority between internal human managers and external agentic-AI platforms. In a two-period model, delegation to AI increases current execution efficiency but reduces retained human recovery capability that becomes valuable when external AI is unavailable or unreliable. The analysis shows that optimal organization is hybrid, that full automation may be suboptimal even when AI dominates task by task, and that platform pricing distorts both delegation and adoption relative to a technological benchmark, especially in longer-horizon execution environments.
\end{abstract}

\noindent\textit{Keywords:} generative AI; authority allocation; organizational design; platform dependence; recovery capability

\noindent\textit{JEL Classification:} D23; L22; L86; O33

\section{Introduction}

Recent advances in generative AI have changed the organizational problem faced by firms. The relevant technology is not merely an internal productivity tool. Increasingly, firms access powerful AI systems through external platforms and use them for multi-step execution rather than one-shot prediction. As a result, the economic issue is not simply whether some tasks can be automated, but how reliance on externally supplied agentic AI reshapes the internal allocation of authority and the preservation of organizational capabilities \citep{AgrawalGansGoldfarb2019,NzembayieUrbano2026}.

In such an environment, the firm must decide which tasks to assign to external AI systems and which tasks to retain under human authority. This question belongs to the classic economics of organizations, where authority, communication, and the distribution of knowledge shape internal decision-making \citep{AghionTirole1997,Dessein2002,Garicano2000}. It is also related to recent work on the allocation of decision authority between humans and artificial intelligence \citep{AtheyBryanGans2020}. However, the organizational problem created by current generative AI differs from earlier formulations in one important respect: the firm delegates not to an internal machine whose use is fully under its control, but to an external platform whose services are priced and whose availability cannot be taken for granted.

This paper develops a simple two-period model of authority allocation with external agentic AI. In period 1, the firm assigns routine tasks to an external AI platform and retains sufficiently exceptional tasks under human authority. In period 2, contingencies may arise in which the external platform becomes unavailable or cannot be fully relied upon, so that the firm must draw on retained internal recovery capability (fallback capability). The model therefore links current authority allocation to future organizational resilience. Relative to the existing literature, I treat AI not as a one-shot signal or decision aid, but as an externally supplied multi-step execution technology, and I allow current delegation choices to affect the firm's future capacity to recover, override, and coordinate when external AI fails or becomes unusable \citep{AgrawalGansGoldfarb2019,AtheyBryanGans2020,NzembayieUrbano2026}.

The analysis yields three main results. First, the optimal organization takes a hybrid form: routine tasks are delegated to AI, while sufficiently exceptional tasks remain under human authority. Second, full automation may be suboptimal even when AI dominates human managers task by task in normal operations. The reason is not that AI necessarily performs poorly on marginal tasks, but that extensive reliance on external AI erodes retained human recovery capability that becomes valuable under contingencies. Third, platform pricing distorts the firm's internal authority structure relative to a technological benchmark, and this distortion is stronger in more agentic environments with longer execution horizons. More broadly, the paper shows that a new externally supplied execution technology can alter not only current task allocation, but also the firm's internal design of control, adaptation, and resilience.

The paper contributes to three strands of literature. It contributes to the economics of authority, communication, and organizational knowledge \citep{AghionTirole1997,Dessein2002,Garicano2000}; to recent work on AI and the allocation of decision authority \citep{AtheyBryanGans2020,AgrawalGansGoldfarb2019}; and to research on technical change, governance, and institutional dependence in platform environments \citep{RochetTirole2003,BakerGibbonsMurphy2023,MarschakWei2024,MukherjeeSaha2024,Kim2025,NzembayieUrbano2026}. In particular, the paper speaks to a JITE-style question: how a new externally supplied technology interacts with internal governance design, adaptation, and organizational control. The central message is that when critical AI capabilities are supplied by external platforms, the firm faces a distinct organizational trade-off between current execution efficiency and the preservation of internal recovery capability.

The remainder of the paper is organized as follows. Section~\ref{sec:lit} discusses the related literature. Section~\ref{sec:model} presents the model. Sections~\ref{sec:baseline}--\ref{sec:capability} derive the main results. Section~\ref{sec:discussion} discusses broader institutional implications, and Section~\ref{sec:conclusion} concludes.

\section{Related Literature}
\label{sec:lit}

This paper first relates to the economics of organizations on authority, communication, and knowledge hierarchies. Formal and real authority have long been central to organizational design \citep{AghionTirole1997}. Communication frictions also shape the efficient allocation of decision rights inside organizations \citep{Dessein2002}. More broadly, the internal organization of knowledge is a core determinant of hierarchy and production \citep{Garicano2000}. The present paper builds on this tradition, but adds a dynamic organizational consideration: authority over current tasks influences the firm's future recovery capability when an important execution technology is externally supplied.

The paper is also closely related to recent work on authority allocation between humans and artificial intelligence. \citet{AtheyBryanGans2020} study how decision authority should be allocated between human and artificial decision-makers. My analysis differs in two respects. First, I model AI not as a one-shot recommendation technology, but as a multi-step execution technology, thereby moving beyond the standard treatment of AI as prediction \citep{AgrawalGansGoldfarb2019}. Second, I emphasize that current delegation decisions affect future internal recovery capability. This channel implies that the case against full automation may persist even when AI dominates human handling task by task in normal operations.

More specifically, the paper is intended to speak to a set of questions that have recently become salient in JITE: how technical change interacts with internal control, how governance adapts to changing organizational environments, and how institutional arrangements shape the use of externally supplied technologies. \citet{MarschakWei2024} show how technical change alters decentralization incentives through the interaction with moral hazard. \citet{MukherjeeSaha2024} analyze how delegation interacts with wage bargaining and innovation. \citet{BakerGibbonsMurphy2023} emphasize that organizational design involves not only incentives, but also control and adaptation through the interaction of formal and relational governance. My paper fits naturally into this broader line of inquiry by asking how a new external execution technology changes the boundary between delegated and retained authority inside the firm.

Finally, the paper connects the literature on technical change and delegation to the economics of platforms and institutional dependence. Because the AI capability considered here is supplied by an external platform, its pricing affects not only whether the technology is adopted, but also how authority is allocated within the firm. In that sense, the paper also speaks to the economics of platforms \citep{RochetTirole2003} and to more recent work on platform-side control and self-preferencing \citep{Kim2025}. The generative-AI setting adds an additional layer of dependence and power asymmetry, which has begun to attract attention in adjacent literatures \citep{NzembayieUrbano2026}. Existing work has not provided a comparably simple theory in which dependence on an external agentic-AI platform jointly shapes current authority allocation, retained human recovery capability, and pricing-induced organizational distortion.

\section{Model}
\label{sec:model}

This section presents a simple two-period model of authority allocation between an external agentic-AI platform and internal human managers. The model is intentionally minimal. Its purpose is not to describe every aspect of generative AI, but to isolate two organizationally central features. First, agentic AI performs \emph{multi-step execution} rather than one-shot prediction. Second, when such AI is supplied by an external platform, current delegation choices affect the firm's future internal recovery capability (fallback capability).

\subsection{Normal Operations}
\label{subsec:normal}

A firm faces a unit mass of tasks indexed by $z\in[0,1]$. Higher $z$ means that a task is more exceptional, more tacit, and harder to verify. The firm chooses a cutoff $s\in[0,1]$: all tasks with $z\le s$ are assigned to an external AI platform, while tasks with $z>s$ remain under human authority.

If task $z$ is assigned to AI, the platform must complete $n\ge 1$ sequential steps. Let the per-step success probability be
\begin{equation}
\label{eq:q}
q(z;\kappa)=\kappa-\alpha z,
\qquad 0<\kappa-\alpha<\kappa<1,
\end{equation}
where $\kappa$ measures AI capability and $\alpha>0$ captures the extent to which performance deteriorates with task exceptionality. Since all $n$ steps must succeed, the overall success probability is
\begin{equation}
\label{eq:Q}
Q(z;\kappa)=q(z;\kappa)^n=(\kappa-\alpha z)^n.
\end{equation}
Equations \eqref{eq:q}--\eqref{eq:Q} provide a reduced-form representation of agentic AI. What matters is not a single prediction, but the probability of completing a sequence of actions without breakdown.

A successfully completed task yields value $v>0$. If AI fails, the firm incurs loss $L>0$. In addition, access to the external AI platform requires a per-step payment $p>0$, so that the total usage cost is $pn$. Hence the expected payoff from assigning task $z$ to AI is
\begin{equation}
\label{eq:uA}
u_A^P(z;\kappa,p)=v-pn-L\bigl(1-Q(z;\kappa)\bigr).
\end{equation}

If task $z$ remains under human authority, expected payoff is
\begin{equation}
\label{eq:uH}
u_H=v-c_H-b,
\end{equation}
where $c_H>0$ is the operating cost of human handling and $b\ge 0$ is a reduced-form agency or bias cost.

The private period-1 gain from assigning task $z$ to AI rather than to humans is therefore
\begin{equation}
\label{eq:DeltaP}
\Delta^P(z;\kappa,p)
\equiv
u_A^P(z;\kappa,p)-u_H
=
c_H+b-pn-L+L(\kappa-\alpha z)^n.
\end{equation}
Because more exceptional tasks are harder for the AI platform, $\Delta^P(z;\kappa,p)$ is decreasing in $z$.

\subsection{Contingencies and Retained Human Recovery Capability}
\label{subsec:contingency}

The organizational reason to retain human authority is not merely that AI may fail on a given task in normal operations. More importantly, heavy reliance on an external platform may erode the firm's internal \emph{recovery capability} (fallback capability), which becomes valuable when the platform cannot be fully relied upon.

I interpret recovery capability as the firm's internally preserved ability to restore operations, override erroneous execution, and coordinate responses when externally supplied AI becomes unavailable or unreliable. To keep the model parsimonious, I represent this capability as an experience stock accumulated through retained human authority in period 1.

Specifically, suppose that the firm's retained human recovery capability after period 1 is
\begin{equation}
\label{eq:H}
H(s)=\bar H+\beta(1-s),
\qquad \beta>0,
\end{equation}
where $\bar H>0$ is baseline capability and $\beta$ measures the amount of recovery-relevant experience preserved by leaving tasks under human authority. Since the mass of tasks retained by humans is $1-s$, \eqref{eq:H} may be interpreted as the aggregate stock of recovery-relevant organizational experience available in period 2.

With probability $\lambda\in(0,1)$, period 2 features a contingency. Conditional on a contingency, the external AI platform remains usable with probability $1-\nu$ and becomes unavailable or unreliable with probability $\nu\in(0,1)$. If the platform remains usable, continuation value is independent of $s$ and can be ignored for the authority-allocation problem. If the platform is unavailable or unreliable, the firm must rely on retained internal recovery capability.

Let
\begin{equation}
\label{eq:rho}
\rho(H(s))
\end{equation}
denote the probability that the firm successfully recovers in that state. Assume that $\rho:\mathbb{R}_+\to[0,1]$ is twice continuously differentiable and satisfies
\begin{equation}
\label{eq:rho-props}
\rho'(H)>0,
\qquad
\rho''(H)<0.
\end{equation}
If recovery succeeds, the firm obtains value $R>0$. Thus, up to an additive constant that does not depend on $s$, the continuation value relevant for the authority choice is
\begin{equation}
\label{eq:continuation}
\lambda \nu R \rho(H(s)).
\end{equation}

This formulation makes the fallback channel more specific than in a generic resilience model. Human authority is valuable not only because humans perform some current tasks, but also because retained authority preserves the internal capability needed to recover when external AI cannot be used or trusted.

\subsection{Private and Technological Objectives}
\label{subsec:objectives}

Let $F\ge 0$ denote the fixed cost of adopting and governing the external AI platform. This fixed cost may be interpreted broadly to include contracting, integration, compliance, and monitoring costs. Relative to the all-human benchmark $s=0$, the firm's private net surplus from choosing cutoff $s$ is
\begin{equation}
\label{eq:SP}
S^P(s;\kappa,p,F)
=
\int_0^s \Delta^P(z;\kappa,p)\,dz
+
\lambda \nu R \bigl[\rho(H(s))-\rho(H(0))\bigr]
-
F\,\mathbf{1}\{s>0\}.
\end{equation}
The indicator term makes clear that $F$ affects whether the platform is adopted at all, but not the internal authority margin conditional on adoption.

Accordingly, conditional on adoption, the firm chooses $s$ to maximize
\begin{equation}
\label{eq:PhiP}
\Phi^P(s;\kappa,p)
=
\int_0^s \Delta^P(z;\kappa,p)\,dz
+
\lambda \nu R \bigl[\rho(H(s))-\rho(H(0))\bigr].
\end{equation}
The first term captures the normal-state gain from delegating a larger set of tasks to AI. The second term captures the loss in future recovery capability when human authority is reduced.

For later comparison, I define a \emph{technological benchmark}. This benchmark does not solve a full social planner's problem. Instead, it asks how the firm would allocate authority if it faced the underlying technological resource cost of AI rather than the private platform price. Let $r\le p$ denote this underlying cost. Then the benchmark gain from assigning task $z$ to AI is
\begin{equation}
\label{eq:DeltaT}
\Delta^T(z;\kappa,r)
=
c_H+b-rn-L+L(\kappa-\alpha z)^n,
\end{equation}
and the corresponding conditional objective is
\begin{equation}
\label{eq:PhiT}
\Phi^T(s;\kappa,r)
=
\int_0^s \Delta^T(z;\kappa,r)\,dz
+
\lambda \nu R \bigl[\rho(H(s))-\rho(H(0))\bigr].
\end{equation}

The distinction between \eqref{eq:PhiP} and \eqref{eq:PhiT} is central. The private problem is governed by the platform's actual access price $p$, whereas the technological benchmark is governed by the underlying resource cost $r$. This distinction allows the analysis to separate two conceptually different questions: whether AI is technologically efficient for the firm to use, and how external platform pricing distorts the firm's internal authority structure.

Finally, adoption occurs if and only if
\begin{equation}
\label{eq:adoption}
\max_{s\in[0,1]} S^P(s;\kappa,p,F)\ge 0.
\end{equation}
Conditional on adoption, however, the internal allocation of authority is determined by the maximization of $\Phi^P(s;\kappa,p)$.

\section{Authority Allocation and Retained Human Recovery Capability}
\label{sec:baseline}

This section derives the model's baseline organizational implications. I first show that the firm's problem has a unique cutoff structure: routine tasks are delegated to the external AI platform, while sufficiently exceptional tasks remain under human authority. I then show that the case against full automation can survive even when AI dominates human handling task by task in normal operations. The reason is intertemporal: authority today shapes retained recovery capability tomorrow.

\subsection{The Optimal Cutoff Structure}
\label{subsec:cutoff}

The firm's conditional objective is given by \eqref{eq:PhiP}. The marginal benefit of increasing the cutoff $s$ is the normal-state gain from assigning the marginal task $s$ to AI rather than to humans, net of the marginal loss in future recovery capability (fallback capability). Because the former falls with task exceptionality while the latter rises as retained capability is depleted, the firm's problem admits a unique threshold.

\begin{proposition}[Unique cutoff structure]
\label{prop:cutoff}
Conditional on adoption, $\Phi^P(s;\kappa,p)$ is strictly concave in $s$. Hence the firm has a unique optimal cutoff $s_P^\ast\in[0,1]$. If
\[
\Delta^P(0;\kappa,p)>\lambda \nu R \beta \rho'(H(0))
\quad\text{and}\quad
\Delta^P(1;\kappa,p)<\lambda \nu R \beta \rho'(H(1)),
\]
then the optimum is interior and characterized by
\begin{equation}
\label{eq:FOC-private-baseline}
\Delta^P(s_P^\ast;\kappa,p)=\lambda \nu R \beta \rho'(H(s_P^\ast)).
\end{equation}
\end{proposition}

The result follows from two monotonicities. First, the normal-state gain from delegating the marginal task to AI is decreasing in task exceptionality, since $\Delta^P(s;\kappa,p)$ falls with $s$. Second, the marginal value of retained human recovery capability rises as the firm automates more tasks, because a higher cutoff reduces $H(s)$ and therefore increases the marginal cost of further capability depletion. These two forces imply a unique crossing and hence a unique organizational cutoff. Appendix~\ref{app:proofs} provides the formal proof.

Proposition~\ref{prop:cutoff} gives the model's basic organizational architecture. The left-hand side of \eqref{eq:FOC-private-baseline} is the period-1 gain from pushing one more marginal task toward AI. The right-hand side is the period-2 value of preserving the marginal unit of human recovery capability. The optimal organization equates these two forces.

\subsection{Why Full Automation May Be Suboptimal}
\label{subsec:no-full-automation}

A static comparison would suggest that if AI outperforms human managers on every task in normal operations, the firm should automate everything. The next result shows why this logic fails in the present setting. When current authority allocation affects future recovery capability, full automation may be suboptimal even if AI dominates human handling task by task.

\begin{proposition}[Why full automation may be suboptimal even when AI dominates task-by-task]
\label{prop:no-full-automation}
Suppose that, in normal operations,
\[
\Delta^P(z;\kappa,p)>0
\qquad
\text{for all } z\in[0,1].
\]
Then:

\begin{enumerate}
\item If
\begin{equation}
\label{eq:no-full-auto-boundary}
\Delta^P(1;\kappa,p)<\lambda \nu R \beta \rho'(\bar H),
\end{equation}
full automation is suboptimal, i.e.
\[
s_P^\ast<1.
\]

\item If, in addition,
\begin{equation}
\label{eq:hybrid-boundary}
\Delta^P(0;\kappa,p)>\lambda \nu R \beta \rho'(H(0)),
\end{equation}
then the unique optimum is hybrid:
\[
0<s_P^\ast<1.
\]

\item The sufficient condition in \eqref{eq:no-full-auto-boundary} becomes easier to satisfy as $\lambda$, $\nu$, $R$, or $\beta$ increase.
\end{enumerate}
\end{proposition}

Part (1) rules out the interpretation that human authority survives only because AI performs poorly on sufficiently exceptional tasks. The assumption $\Delta^P(z;\kappa,p)>0$ excludes that possibility. Instead, the result follows because the marginal organizational cost of full automation is
\[
\lambda \nu R \beta \rho'(\bar H),
\]
which is increasing in contingency risk, platform-unusability risk, recovery value, and the rate at which authority retention preserves capability. Part (2) then identifies a parameter region in which the firm chooses a genuinely hybrid authority structure. Appendix~\ref{app:proofs} provides the formal proof.

The economic content of Proposition~\ref{prop:no-full-automation} is therefore sharper than in a purely static automation model. Even when AI dominates task by task in normal operations, broad delegation can still be suboptimal because authority today determines how much internal recovery capability survives into abnormal states tomorrow.

\section{Platform Pricing and Organizational Distortion}
\label{sec:pricing}

The previous section showed that the firm may rationally preserve human authority in order to protect future recovery capability. This section shows that platform pricing does more than affect whether AI is adopted. It also distorts the firm's internal allocation of authority relative to a technological benchmark, and this distortion is stronger in more agentic environments with longer execution horizons.

The key observation is simple. The benchmark problem differs from the firm's private problem only in the price of access to AI. In the private problem, the firm pays the platform price $p$ per execution step. In the technological benchmark, the relevant cost is the underlying resource cost $r\le p$. Since the pricing wedge $p-r$ applies step by step, it lowers the marginal return to authority reallocation toward AI even when the underlying technology is unchanged.

\subsection{The Distortion in Internal Authority Allocation}
\label{subsec:pricing-distortion}

Recall that conditional on adoption, the firm's private cutoff $s_P^\ast$ solves
\[
\Delta^P(s_P^\ast;\kappa,p)=\lambda \nu R \beta \rho'(H(s_P^\ast)).
\]
The technological benchmark chooses $s_T^\ast$ to satisfy
\begin{equation}
\label{eq:FOC-tech-pricing}
\Delta^T(s_T^\ast;\kappa,r)=\lambda \nu R \beta \rho'(H(s_T^\ast)).
\end{equation}
The continuation term is the same in both problems. The only difference is the per-step pricing wedge embedded in $\Delta^P$ relative to $\Delta^T$.

\begin{proposition}[Platform pricing distorts the internal allocation of authority]
\label{prop:pricing-distortion}
Suppose $p>r$ and that both the private problem and the technological benchmark admit interior solutions. Then
\[
s_P^\ast<s_T^\ast.
\]

Moreover, for every task $z\in[0,1]$,
\[
\Delta^T(z;\kappa,r)-\Delta^P(z;\kappa,p)=n(p-r),
\]
and for every delegated mass $s>0$,
\[
\Phi^T(s;\kappa,r)-\Phi^P(s;\kappa,p)=sn(p-r).
\]
Hence the pricing-induced distortion is larger in more agentic environments with longer execution horizons.
\end{proposition}

The logic is immediate. The pricing wedge shifts down the private marginal gain from delegation by exactly $n(p-r)$ for every task. Since the continuation term is identical across the two problems and the marginal gain from delegation is decreasing in task exceptionality, the benchmark zero crossing must lie to the right of the private one. The same wedge also scales linearly with the execution horizon $n$, so the distortion is stronger precisely when AI use is more multi-step and hence more agentic. Appendix~\ref{app:proofs} gives the formal derivation.

Proposition~\ref{prop:pricing-distortion} is important because it identifies an \emph{organizational} distortion rather than a purely financial one. The platform's pricing policy does not simply alter the firm's total cost of using AI. It changes the firm's internal authority architecture by shifting tasks that would be technologically efficient to delegate back toward human authority.

\subsection{A Remark on the Adoption Margin}
\label{subsec:adoption-margin}

The pricing wedge also affects the extensive margin of adoption. From \eqref{eq:SP}, adoption occurs if and only if
\[
\max_{s\in[0,1]} \Phi^P(s;\kappa,p)\ge F.
\]
Because the benchmark enjoys a uniformly higher gain from delegation, there exist parameter regions in which AI adoption is technologically efficient but privately unattractive.

\begin{remark}
\label{rem:adoption-margin}
If $p>r$, then for every $s>0$,
\[
\Phi^T(s;\kappa,r)-\Phi^P(s;\kappa,p)=sn(p-r)>0.
\]
Therefore, whenever
\[
\max_{s\in[0,1]} \Phi^T(s;\kappa,r)\ge F>\max_{s\in[0,1]} \Phi^P(s;\kappa,p),
\]
the technological benchmark adopts AI while the private firm does not.
\end{remark}

Remark~\ref{rem:adoption-margin} complements Proposition~\ref{prop:pricing-distortion}. The proposition concerns the intensive margin: conditional on adoption, platform pricing leads the firm to delegate too little. The remark concerns the extensive margin: for sufficiently large fixed costs, platform pricing may prevent adoption altogether even when adoption would be efficient under the technological benchmark.

\section{AI Capability and the Persistence of Human Authority}
\label{sec:capability}

The previous sections showed that the firm may rationally preserve human authority in order to protect future recovery capability, and that platform pricing distorts this choice relative to a technological benchmark. A natural question is whether continued improvement in AI capability eventually eliminates human authority altogether. This section shows that the answer is generally no. Better AI expands delegation, but it need not drive the optimal domain of human authority to zero.

\subsection{Better AI and the Expansion of Delegation}
\label{subsec:better-ai}

An increase in $\kappa$ raises the success probability of AI on every task and therefore increases the normal-state gain from delegation. It is therefore natural to expect the firm to assign a broader set of tasks to the external platform. The next result establishes this comparative static and then sharpens it by showing that, under a transparent sufficient condition, a strictly positive domain of human authority survives uniformly over all admissible AI capability levels.

\begin{proposition}[Better AI expands delegation, but positive human authority may persist uniformly]
\label{prop:better-ai-not-full}
Assume that the private problem admits an interior solution $s_P^\ast(\kappa)\in(0,1)$ for $\kappa$ in an interval $\mathcal{K}\subset(\alpha,1]$. Then:

\begin{enumerate}
\item $s_P^\ast(\kappa)$ is strictly increasing in $\kappa$.

\item Suppose there exists $\varepsilon\in(0,1)$ such that
\begin{equation}
\label{eq:uniform-persistence-condition}
\Delta^P(1-\varepsilon;1,p)
<
\lambda \nu R \beta \rho'(\bar H+\beta\varepsilon).
\end{equation}
Then
\[
s_P^\ast(\kappa)< 1-\varepsilon
\qquad
\text{for all } \kappa\in\mathcal{K}.
\]
Thus, a strictly positive share greater than $\varepsilon$ of tasks remains under human authority uniformly across all admissible AI capability levels.

\item The sufficient condition in \eqref{eq:uniform-persistence-condition} becomes easier to satisfy as $\lambda$, $\nu$, or $R$ increase.
\end{enumerate}
\end{proposition}

Part (1) is the natural comparative static: better AI expands the range of tasks delegated to the platform. Part (2) is stronger than the claim that full automation may fail. It shows that, under an economically interpretable sufficient condition, AI progress still leaves a uniformly positive domain of human authority. Part (3) clarifies when this is most likely to occur: precisely when contingencies are more important, external AI is more likely to be unusable in those contingencies, or successful recovery is more valuable. Appendix~\ref{app:proofs} provides the formal derivation.

The intuition is straightforward. A higher $\kappa$ improves current execution and therefore pushes the firm toward more delegation. But that force operates only on the period-1 margin. It does not remove the period-2 value of retained internal recovery capability. As long as contingency-side value remains sufficiently large, the firm preserves a positive amount of human authority in order to sustain the capacity to intervene, override, and recover when the external platform cannot be fully relied upon.

This result sharpens the economic meaning of the model. The relevant organizational question is not whether AI becomes ``good enough'' in a purely technical sense. Rather, it is whether improvements in current execution dominate the future value of internal recovery capability. Proposition~\ref{prop:better-ai-not-full} shows that the answer can be negative even in a strong sense: not only may full automation fail, but a strictly positive human authority share may survive uniformly over all admissible AI capability levels.

Taken together, Propositions~\ref{prop:no-full-automation}, \ref{prop:pricing-distortion}, and \ref{prop:better-ai-not-full} imply that dependence on external agentic-AI platforms creates a persistent role for human authority inside the firm. Platform pricing distorts the private extent of delegation, and even absent such distortion, the firm may rationally preserve human control as an organizational buffer against contingencies.

\section{Discussion}
\label{sec:discussion}

The model developed in this paper is deliberately simple, but its implications are broader than a narrow automation problem. The central point is that when agentic AI is supplied by an external platform, the firm's organizational problem is not exhausted by current execution efficiency. Delegation to external AI changes the firm's internal authority structure because it also changes the stock of retained human recovery capability available when contingencies arise. In that sense, the relevant trade-off is not merely between human and machine performance on current tasks, but between present efficiency and future organizational resilience.

This perspective helps clarify why full automation need not be the natural limiting outcome of AI progress. In a purely static model, sufficiently capable AI would eventually displace human authority wherever it outperforms human handling. In the present framework, however, the value of human authority is not confined to contemporaneous task execution. Human authority also preserves internal recovery capability for intervention, override, and coordination under abnormal conditions, especially when externally supplied AI becomes unavailable or cannot be fully relied upon. The model therefore suggests that the persistence of human authority need not reflect technological backwardness. It may instead be an optimal organizational response to dependence on an external execution technology.

A related implication concerns the institutional role of platform pricing. In the model, pricing does not simply influence whether firms adopt AI. It also affects which decisions remain internal to the organization. This is important because it implies that the governance of external AI platforms has consequences inside the firm. A priced access regime changes the boundary between delegated and retained tasks, and therefore changes the scope of human discretion, oversight, and organizational learning. Because the pricing wedge is paid step by step, these distortions are stronger in settings where AI use is more genuinely agentic, that is, more reliant on long-horizon execution.

The model's notion of retained human recovery capability should be interpreted broadly as an internal experience stock relevant for abnormal states. It may represent crisis-response ability, override capacity, supervisory judgment, tacit operational knowledge, or the practical competence needed to diagnose and correct failures when the external platform cannot be fully relied upon. The reduced-form specification is intentionally parsimonious, but the underlying idea is general: internal authority sustains a stock of organizational capability that is valuable precisely when routine external execution becomes insufficient.

Several limitations are worth noting. First, the model takes platform pricing as exogenous. This is appropriate for isolating the firm's internal organization problem, but it leaves aside the upstream determination of access terms and the market structure of AI provision. Second, the model represents contingencies in reduced form rather than deriving them from a richer microfoundation of breakdown, override, or domain shift. Third, the firm is represented as a single decision-making unit, abstracting from strategic interaction across firms. These simplifications are deliberate, but they also indicate natural directions for future work.

Two extensions seem especially promising. One is to endogenize platform pricing or platform market power, thereby linking internal authority allocation to the industrial organization of AI provision. The other is to connect the present mechanism to make-or-buy decisions and firm boundaries. If external AI platforms perform routine execution while internal teams preserve recovery and exception-handling capability, then the trade-off identified here may also shape outsourcing, vertical scope, and the internal structure of the firm. These extensions lie beyond the present paper, but the baseline mechanism developed here provides a tractable starting point for them.

\section{Conclusion}
\label{sec:conclusion}

This paper has studied authority allocation in a firm that relies on an external agentic-AI platform for task execution. The central argument is that the organizational consequences of generative AI cannot be understood by looking only at current task-level efficiency. Because delegation to external AI affects the firm's retained human recovery capability, the relevant organizational trade-off is intertemporal: greater current execution efficiency may come at the cost of weaker future resilience.

The analysis delivers three main conclusions. First, the optimal organization generally takes a hybrid form in which routine tasks are delegated to AI while sufficiently exceptional tasks remain under human authority. Second, full automation may be suboptimal even when AI dominates human handling task by task in normal operations, because authority today helps preserve internal recovery capability for future contingencies. Third, platform pricing distorts the firm's internal authority structure relative to a technological benchmark, affecting not only whether AI is adopted but also how authority is allocated conditional on adoption, with stronger distortions in longer-horizon execution environments.

More broadly, the paper suggests that the economics of generative AI is not only about the substitution of machines for labor. It is also about how dependence on externally supplied AI reshapes internal governance, organizational learning, and the boundaries of human discretion. In that sense, the relevant question is not simply how much AI a firm should use, but how a firm should organize itself when critical execution capabilities are supplied from outside the organization.


\appendix
\section{Proofs and Additional Derivations}
\label{app:proofs}

\subsection*{Proof of Proposition~\ref{prop:cutoff}}

Recall that the firm's conditional objective is
\[
\Phi^P(s;\kappa,p)
=
\int_0^s \Delta^P(z;\kappa,p)\,dz
+
\lambda \nu R \bigl[\rho(H(s))-\rho(H(0))\bigr].
\]
Differentiating with respect to $s$ yields
\begin{equation}
\label{eq:PhiP-prime-app}
\frac{\partial \Phi^P(s;\kappa,p)}{\partial s}
=
\Delta^P(s;\kappa,p)-\lambda \nu R \beta \rho'(H(s)).
\end{equation}
Using \eqref{eq:DeltaP} and \eqref{eq:H},
\[
\frac{\partial \Delta^P(s;\kappa,p)}{\partial s}
=
-Ln\alpha(\kappa-\alpha s)^{n-1}<0,
\]
because $L>0$, $n\ge 1$, $\alpha>0$, and $\kappa-\alpha s>0$ on $[0,1]$ by \eqref{eq:q}. Moreover,
\[
\frac{d}{ds}\!\left[-\lambda \nu R \beta \rho'(H(s))\right]
=
-\lambda \nu R \beta \rho''(H(s))H'(s)
=
\lambda \nu R \beta^2 \rho''(H(s))<0,
\]
since $H'(s)=-\beta$ and $\rho''(H)<0$ by \eqref{eq:rho-props}. Hence
\begin{equation}
\label{eq:PhiP-second-app}
\frac{\partial^2 \Phi^P(s;\kappa,p)}{\partial s^2}
=
-Ln\alpha(\kappa-\alpha s)^{n-1}
+
\lambda \nu R \beta^2 \rho''(H(s))
<0.
\end{equation}
Therefore $\Phi^P$ is strictly concave on $[0,1]$ and admits a unique maximizer $s_P^\ast$.

If
\[
\Delta^P(0;\kappa,p)>\lambda \nu R \beta \rho'(H(0))
\quad\text{and}\quad
\Delta^P(1;\kappa,p)<\lambda \nu R \beta \rho'(H(1)),
\]
then \eqref{eq:PhiP-prime-app} is positive at $s=0$ and negative at $s=1$. Since \eqref{eq:PhiP-prime-app} is strictly decreasing by \eqref{eq:PhiP-second-app}, it crosses zero exactly once in $(0,1)$. Hence the unique maximizer is interior and satisfies
\[
\Delta^P(s_P^\ast;\kappa,p)=\lambda \nu R \beta \rho'(H(s_P^\ast)).
\]
This proves Proposition~\ref{prop:cutoff}.

\subsection*{Proof of Proposition~\ref{prop:no-full-automation}}

Let
\[
G^P(s)\equiv \Delta^P(s;\kappa,p)-\lambda \nu R \beta \rho'(H(s)).
\]
By the same argument used in the proof of Proposition~\ref{prop:cutoff}, the objective $\Phi^P(s;\kappa,p)$ is strictly concave, and
\[
\frac{d\Phi^P(s;\kappa,p)}{ds}=G^P(s).
\]
Moreover,
\[
\frac{dG^P(s)}{ds}
=
-Ln\alpha(\kappa-\alpha s)^{n-1}
+
\lambda \nu R \beta^2 \rho''(H(s))
<0,
\]
because $\rho''(H(s))<0$. Hence $G^P(s)$ is strictly decreasing in $s$.

To prove part (1), evaluate the derivative at full automation:
\[
G^P(1)=\Delta^P(1;\kappa,p)-\lambda \nu R \beta \rho'(H(1)).
\]
Since $H(1)=\bar H$, condition \eqref{eq:no-full-auto-boundary} is exactly
\[
G^P(1)<0.
\]
Because $\Phi^P$ is strictly concave, a negative derivative at the right boundary implies that the optimum cannot occur at $s=1$. Hence
\[
s_P^\ast<1.
\]

To prove part (2), condition \eqref{eq:hybrid-boundary} implies
\[
G^P(0)=\Delta^P(0;\kappa,p)-\lambda \nu R \beta \rho'(H(0))>0.
\]
Together with $G^P(1)<0$ from part (1), strict monotonicity of $G^P$ implies that there exists a unique zero in $(0,1)$. By strict concavity, that zero is the unique maximizer, so
\[
0<s_P^\ast<1.
\]

For part (3), the right-hand side of \eqref{eq:no-full-auto-boundary} is
\[
\lambda \nu R \beta \rho'(\bar H),
\]
which is strictly increasing in $\lambda$, $\nu$, $R$, and $\beta$, since $\rho'(\bar H)>0$. Hence the sufficient condition ruling out full automation becomes easier to satisfy as any of these parameters increases. This proves Proposition~\ref{prop:no-full-automation}.

\subsection*{Proof of Proposition~\ref{prop:pricing-distortion}}

Recall
\[
\Delta^P(z;\kappa,p)
=
c_H+b-pn-L+L(\kappa-\alpha z)^n
\]
and
\[
\Delta^T(z;\kappa,r)
=
c_H+b-rn-L+L(\kappa-\alpha z)^n.
\]
Hence, for every $z\in[0,1]$,
\begin{equation}
\label{eq:Delta-gap-app}
\Delta^T(z;\kappa,r)-\Delta^P(z;\kappa,p)=n(p-r).
\end{equation}
If $p>r$, then \eqref{eq:Delta-gap-app} is strictly positive and increasing linearly in $n$.

Define
\[
G^P(s)\equiv \Delta^P(s;\kappa,p)-\lambda \nu R \beta \rho'(H(s))
\]
and
\[
G^T(s)\equiv \Delta^T(s;\kappa,r)-\lambda \nu R \beta \rho'(H(s)).
\]
Interior solutions satisfy
\[
G^P(s_P^\ast)=0
\qquad\text{and}\qquad
G^T(s_T^\ast)=0.
\]
By \eqref{eq:Delta-gap-app},
\begin{equation}
\label{eq:G-gap-app}
G^T(s)-G^P(s)=n(p-r)>0
\qquad
\text{for all } s\in[0,1].
\end{equation}

Next, both $G^P$ and $G^T$ are strictly decreasing in $s$. Indeed,
\[
\frac{d\Delta^P(s;\kappa,p)}{ds}
=
-Ln\alpha(\kappa-\alpha s)^{n-1}<0,
\]
\[
\frac{d\Delta^T(s;\kappa,r)}{ds}
=
-Ln\alpha(\kappa-\alpha s)^{n-1}<0,
\]
and
\[
\frac{d}{ds}\!\left[-\lambda \nu R \beta \rho'(H(s))\right]
=
\lambda \nu R \beta^2 \rho''(H(s))<0.
\]
Therefore
\[
\frac{dG^P(s)}{ds}<0
\qquad\text{and}\qquad
\frac{dG^T(s)}{ds}<0.
\]

Since $G^T$ is a strict upward shift of $G^P$ by \eqref{eq:G-gap-app}, while both functions are strictly decreasing, the unique zero of $G^T$ must lie strictly to the right of the unique zero of $G^P$. Hence
\[
s_P^\ast<s_T^\ast.
\]

It remains to show the total-objective gap. Using \eqref{eq:PhiP} and \eqref{eq:PhiT},
\[
\Phi^T(s;\kappa,r)-\Phi^P(s;\kappa,p)
=
\int_0^s \bigl[\Delta^T(z;\kappa,r)-\Delta^P(z;\kappa,p)\bigr]dz,
\]
because the continuation term is identical in both problems. Substituting \eqref{eq:Delta-gap-app},
\[
\Phi^T(s;\kappa,r)-\Phi^P(s;\kappa,p)
=
\int_0^s n(p-r)dz
=
sn(p-r).
\]
Thus the pricing-induced distortion grows linearly with the execution horizon $n$. This proves Proposition~\ref{prop:pricing-distortion}.

\subsection*{Derivation for Remark~\ref{rem:adoption-margin}}

If $p>r$, then by the final line of the proof above,
\[
\Phi^T(s;\kappa,r)-\Phi^P(s;\kappa,p)=sn(p-r)>0
\qquad
\text{for all } s>0.
\]
Adoption in the private problem occurs if and only if
\[
\max_{s\in[0,1]} \Phi^P(s;\kappa,p)\ge F.
\]
Therefore, if
\[
\max_{s\in[0,1]} \Phi^T(s;\kappa,r)\ge F>\max_{s\in[0,1]} \Phi^P(s;\kappa,p),
\]
the technological benchmark adopts while the private firm does not. This establishes Remark~\ref{rem:adoption-margin}.

\subsection*{Proof of Proposition~\ref{prop:better-ai-not-full}}

For interior solutions, define
\[
G(s,\kappa)\equiv \Delta^P(s;\kappa,p)-\lambda \nu R \beta \rho'(H(s)).
\]
By Proposition~\ref{prop:cutoff}, the optimal cutoff $s_P^\ast(\kappa)$ is the unique solution to
\begin{equation}
\label{eq:G-zero-app}
G(s_P^\ast(\kappa),\kappa)=0.
\end{equation}

To prove part (1), differentiate $G$ with respect to $\kappa$. Using \eqref{eq:DeltaP},
\[
\frac{\partial G(s,\kappa)}{\partial \kappa}
=
\frac{\partial \Delta^P(s;\kappa,p)}{\partial \kappa}
=
Ln(\kappa-\alpha s)^{n-1}>0.
\]
Next, differentiate $G$ with respect to $s$:
\[
\frac{\partial G(s,\kappa)}{\partial s}
=
\frac{\partial \Delta^P(s;\kappa,p)}{\partial s}
-\lambda \nu R \beta \frac{d}{ds}\rho'(H(s)).
\]
Since
\[
\frac{\partial \Delta^P(s;\kappa,p)}{\partial s}
=
-Ln\alpha(\kappa-\alpha s)^{n-1}<0
\]
and
\[
-\lambda \nu R \beta \frac{d}{ds}\rho'(H(s))
=
-\lambda \nu R \beta \rho''(H(s))H'(s)
=
\lambda \nu R \beta^2 \rho''(H(s))<0,
\]
it follows that
\[
\frac{\partial G(s,\kappa)}{\partial s}<0.
\]
Therefore, by the implicit function theorem applied to \eqref{eq:G-zero-app},
\[
\frac{d s_P^\ast(\kappa)}{d\kappa}
=
-
\frac{\partial G/\partial \kappa}{\partial G/\partial s}
>0.
\]
This proves part (1).

For part (2), suppose there exists $\varepsilon\in(0,1)$ such that
\[
\Delta^P(1-\varepsilon;1,p)
<
\lambda \nu R \beta \rho'(\bar H+\beta\varepsilon).
\]
Since $H(1-\varepsilon)=\bar H+\beta\varepsilon$, this is equivalent to
\[
G(1-\varepsilon,1)<0.
\]
Now $\Delta^P(s;\kappa,p)$ is increasing in $\kappa$, so for every $\kappa\in\mathcal{K}\subset(\alpha,1]$,
\[
\Delta^P(1-\varepsilon;\kappa,p)\le \Delta^P(1-\varepsilon;1,p).
\]
Therefore
\[
G(1-\varepsilon,\kappa)\le G(1-\varepsilon,1)<0
\qquad
\text{for all } \kappa\in\mathcal{K}.
\]
Because $G(s,\kappa)$ is strictly decreasing in $s$ and has a unique zero at $s=s_P^\ast(\kappa)$, the inequality
\[
G(1-\varepsilon,\kappa)<0
\]
implies
\[
s_P^\ast(\kappa)< 1-\varepsilon
\qquad
\text{for all } \kappa\in\mathcal{K}.
\]
Hence a strictly positive share greater than $\varepsilon$ of tasks remains under human authority uniformly over all admissible AI capability levels.

For part (3), the right-hand side of \eqref{eq:uniform-persistence-condition} is
\[
\lambda \nu R \beta \rho'(\bar H+\beta\varepsilon),
\]
which is strictly increasing in $\lambda$, $\nu$, and $R$, since $\beta>0$ and $\rho'(\bar H+\beta\varepsilon)>0$. Hence the sufficient condition becomes easier to satisfy as any of these parameters increases. This proves Proposition~\ref{prop:better-ai-not-full}.



\printbibliography

\end{document}
